% Katalog H: Opt-in-Module (Palette Slot 8).
\ZSFNewColumn
\StartChapter[kat:module]{H · Module}

\begin{runintext}
  Fünf Module, die eine ZSF einzeln in \ZSFlabel{preamble.tex} zuschaltet.
  Dieser Katalog lädt alle, damit sie beurteilbar sind — ein Fach-Fork
  startet ohne sie.
\end{runintext}

\SubsectionBar[kat:mod-math]{11\_math\_advanced}
\ZSFindex{Operator}

\begin{formulabox}[H-01 · Erweiterte Operatoren]
  \[
    \begin{aligned}
      &\Ker A,\quad \rang A,\quad \Spur A,\quad \diag(A),\\
      &\spanop(v_1,v_2),\quad \eig(A),\quad \proj_U(v),\\
      &\sig(A),\quad \Adj(A),\quad \Re z,\quad \Im z,\quad
        \asterisk,\quad \oasterisk
    \end{aligned}
  \]
\end{formulabox}
\Zweck{Für LinAlg und Analysis. Ohne das Modul sind diese Namen nicht
definiert.}

\begin{formulabox}[H-02 · Klammer über Zellgrenzen]
  \[
    A=\begin{pmatrix}
      1 & 0 & 4 \\
      \tikzmark[xshift=-\ZSFspaceS]{kat-brace-left}0 &
      1\tikzmark[xshift=\ZSFspaceS]{kat-brace-right} & 5
    \end{pmatrix}
  \]
  \drawbrace[brace mirrored]{kat-brace-left}{kat-brace-right}
  \annote[yshift=-\ZSFspaceM]{brace-\thebrace}{\ZSFfontDiagramLabel Mustertext}
  \ZSFbraceRoom
\end{formulabox}
\Zweck{Der Fall, den ZSFbraceunder nicht kann. Zeichnet als Overlay ohne
eigene Höhe — den Platz stellt ZSFbraceRoom.}

\SubsectionBar[kat:mod-plots]{12\_plots}
\ZSFindex{Plot}

\begin{defbox}[H-03 · pgfplots][weight=quiet]
  Lädt pgfplots für die \ZSFlabel{axis}-Umgebung. Vorgeführt in
  \ZSFref{kat:bilder}, weil ein Plot dort in seinem Container steht.
\end{defbox}
\Zweck{Nur für Dokumente mit echten Plot-Inhalten; tikzpicture allein
braucht das Modul nicht.}

\SubsectionBar[kat:mod-code]{65\_code\_style}
\ZSFindex{codebox}
\ZSFindex{part}

\begin{codebox}[H-04 · codebox]
\begin{lstlisting}[style=CodeExpert]
# Mustertext zum Formvergleich
def beispiel(text="Zeichenkette", zahl=42):
    if zahl > 0:
        return text
    return None
\end{lstlisting}
\end{codebox}
\Zweck{Code-Container mit CodeExpert-Stil; Kommentare sind nativ, nicht über
CodeLine.}

\begin{codebox}[H-05 · part=first][part=first]
\begin{lstlisting}[style=CodeExpert]
# Rahmen oben, unten offen
werte = [1, 2, 3]
\end{lstlisting}
\end{codebox}
\begin{codebox}[][part=mid]
\begin{lstlisting}[style=CodeExpert]
# part=mid: beide Kanten offen
summe = sum(werte)
\end{lstlisting}
\end{codebox}
\begin{codebox}[][part=last]
\begin{lstlisting}[style=CodeExpert]
# part=last: Rahmen unten
print(summe)
\end{lstlisting}
\end{codebox}
\Zweck{Mehrteilige Gruppe mit offenen Kanten. Vorbelegung: whole.}

\SubsectionBar[kat:mod-comments]{67\_code\_comments}
\ZSFindex{CodeLine}

\begin{defbox}[H-06 · CodeLine][weight=quiet]
  \ttfamily
  \CodeLine{kurz = 1}[passt rechts daneben]\\
  \CodeLine{sehr_langer_variablenname = 1}[zu breit, deshalb darüber]\\
  \CodeLine{limit = 10}\InlineComment{H-07 · InlineComment}\\
  \OverlineComment{H-08 · OverlineComment}
  \CodeLine{fertig = True}
\end{defbox}
\Zweck{Entscheidet nach Breite, wo der Kommentar steht. Wirkt in jeder
Textbox, aber nicht in lstlisting — unabhängig von 65.}

\SubsectionBar[kat:mod-index]{66\_index}
\ZSFindex{Stichwortverzeichnis}
\ZSFindexsee{Register}{Stichwortverzeichnis}

\begin{defbox}[H-09 · ZSFindex][weight=quiet]
  Unsichtbarer Registereintrag für Begriffe, die nicht als
  \ZSFkeyword*{ZSFkeyword} im Text stehen. Jeder Eintrag dieses Katalogs
  setzt einen — das Ergebnis steht am Dokumentende.
\end{defbox}
\Zweck{Für Box-Titel, Verfahrensnamen und Tabellen-Themen.}

\begin{defbox}[H-10 · ZSFindexsee][weight=quiet]
  Verweis von einem Synonym auf den kanonischen Begriff, samt farbcodierter
  Zielnummer. Im Register steht »Register, siehe Stichwortverzeichnis«.
\end{defbox}
\Zweck{Abkürzungen, Eponyme, Einheiten und Symbole auf den Sachbegriff
führen — der wertvollste Hebel für die Auffindbarkeit.}
